\section{中都之后：北方治理不是一张总账}

1136--1208年的金朝在位纪年，跨越刘豫政权废除后的直接经营、1141年绍兴和议、1153年迁中都、1161年南侵受挫、1164年隆兴和议，以及1206--1208年宋金战事与嘉定和议。连续纪年提示制度并未随一次战争停止，但也不允许我们把华北、河南、东北和宋金边区看成同样被治理的区域。

《金史》的百官、食货与地理材料能给出中央制度和术语，金代碑志、城址、窑冶与钱币遗存则能补问迁徙、城市生产和区域交通。二者不能互相替代：中都的城建或一处冶铸遗址不代表所有路，墓志的官衔还须区分实授、追赠和检校。金朝的“北方统治”应当被理解为不同地方在道路、军镇、税源和官员网络中被接入的程度，而不是民族标签的简单对立。

\section{边界所需的财政与地方社会}

1136年以后，边界的日常成本不只来自大战。驻军、转运、屯田、桥路、市场和地方征收共同决定和议能否维持。1141年与1164年的协议可以证明外交框架，不能证明边境社会从此安定；1250年以前的每一笔军费或贡赐也不能由正史总数直接换算为家庭负担。讨论金的财政，应分别追问记录的机构、地域、币制和报告情境。

1153年中都成为都城，1161年渡江受挫，这两件事尤其说明城市、漕运和军事行动之间的关系。迁都把中枢更深地放入华北的道路和仓储网络，南侵却显示长距离供给与河流战场的限度。其间的工匠、商人、军户和被迁徙者并不只以“国家资源”存在；可见他们的往往是局部碑刻、契约或遗址，必须抵抗用一个城市或一条法令概括整个金朝的冲动。
